% =====================================================================
% R2.7 — "The cingulate carries different content in different rhythms"
% Seventh / closing paragraph of Result 2 (flagship). Synthesizes R2.5 (low-gamma
% cingulate MEMORY) + R2.6 (beta cingulate INFERENCE) into the multiplexing twist.
%
% ESTIMATOR = rho_sym throughout. NO new compute; reads R2.5 (audit_159) + R2.6 (audit_160).
% THE 2x2 (rho_sym, cingulate; verified 2026-07-07):
%              low_gamma            beta
%   ENCODING   q=0.035 CLEARS       p=0.235/0.073 null   (audit_159 / audit_158)
%   INFERENCE  p=0.93   null        q=0.030/0.040 CLEARS (audit_159 / audit_160)
%   -> cingulate carries MEMORY at low-gamma and INFERENCE at beta; the other
%      component is null in each rhythm. One region, content switched by frequency.
% STRENGTH OF EACH HALF (state honestly): memory half = cohort-clearing absolute trace
%   (low-gamma, R2.5, solid). Inference half = the LEADING CANDIDATE of R2.6 (clears the
%   full-length strength null; length-matched control pending) -> provisional. So the
%   multiplexing reading is FIRM for memory, PROVISIONAL for inference. Outline: R2.7
%   "*(Suggestive - the inference half is a hint.)*"
% Guardrails: state the asymmetric confidence (firm memory / provisional inference); no
%   behavioural claim; no absence-narration; register = NN results (feedback_results_not_story).
% Provenance: audit_159 + audit_160 (rho_sym); headline 02_encoding_vs_inference.md (N2, "honest twist");
%   outline R2.7.
% Prose unwrapped (one line per paragraph); American English.
% =====================================================================

\paragraph{The cingulate carries different content in different rhythms.}
The two cingulate results describe one region holding different task content at different frequencies. In \(\lowgamma\) the cingulate carries the memory trace and not the inference component; in \(\beta\) it carries the inference component and not the memory trace. The premises the patient was shown surface there in one rhythm and the relations the patient reasoned out in another --- the same cortex, read at two frequencies, holding two different parts of the task. The two halves stand on different footing: the memory reading rests on a cohort-clearing trace, while the inference reading is the directional lead of the previous paragraph --- consistent in sign but length-assisted --- so the picture is firm for memory and provisional for inference.
